\documentclass{article}

% Packages
\usepackage{fancyhdr}
\usepackage{extramarks}
\usepackage{amsmath}
\usepackage{amsthm}
\usepackage{amsfonts}
\usepackage{tikz}
\usepackage[plain]{algorithm}
\usepackage{algpseudocode}
\usepackage{enumerate}
\usepackage{dsfont}
\usepackage{bbm}

\usetikzlibrary{automata,positioning}

% Document Layout
\topmargin=-0.45in
\evensidemargin=0in
\oddsidemargin=0in
\textwidth=6.5in
\textheight=9.0in
\headsep=0.25in
\linespread{1.1}

% Page Style
\pagestyle{fancy}
\lhead{\hmwkAuthorName}
\chead{\hmwkClass:\ \hmwkTitle}
\rhead{Section \hmwkSection, \firstxmark}
\lfoot{\lastxmark}
\cfoot{\thepage}
\renewcommand\headrulewidth{0.4pt}
\renewcommand\footrulewidth{0.4pt}

% Paragraph Settings
\setlength\parindent{0pt}
\setlength{\parskip}{5pt}

% Section Management
\newcommand{\hmwkSection}{A} % Current section (A, B, or C) - update manually

% Problem Header Management
\newcommand{\enterProblemHeader}[1]{
  \nobreak\extramarks{}{Problem \arabic{#1} continued on next page\ldots}\nobreak{}
  \nobreak\extramarks{Problem \arabic{#1} (continued)}{Problem \arabic{#1} continued on next page\ldots}\nobreak{}
}

\newcommand{\exitProblemHeader}[1]{
  \nobreak\extramarks{Problem \arabic{#1} (continued)}{Problem \arabic{#1} continued on next page\ldots}\nobreak{}
  \stepcounter{#1}
  \nobreak\extramarks{Problem \arabic{#1}}{}\nobreak{}
}

% Counters
\setcounter{secnumdepth}{0}
\newcounter{partCounter}
\newcounter{homeworkProblemCounter}
\setcounter{homeworkProblemCounter}{1}
\nobreak\extramarks{Problem \arabic{homeworkProblemCounter}}{}\nobreak{}

% Homework Problem Environment
% Optional argument adjusts problem counter for non-sequential problems
\newenvironment{homeworkProblem}[1][-1]{
  \ifnum#1>0
    \setcounter{homeworkProblemCounter}{#1}
  \fi
  \section{Problem \arabic{homeworkProblemCounter}}
  \setcounter{partCounter}{1}
  \enterProblemHeader{homeworkProblemCounter}
}{
  \exitProblemHeader{homeworkProblemCounter}
}

% Assignment Details
\newcommand{\hmwkTitle}{Sheet\ \#3}
\newcommand{\hmwkDueDate}{Mar 11, 2026}
\newcommand{\hmwkClass}{C6.2 Continuous Optimization}
\newcommand{\hmwkClassInstructor}{Professor C. Cartis}
\newcommand{\hmwkAuthorName}{\textbf{Ray Tsai}}

% Title Page
\title{
  \vspace{2in}
  \textmd{\textbf{\hmwkClass:\ \hmwkTitle}}\\
  \normalsize\vspace{0.1in}\small{Due\ on\ \hmwkDueDate\ at 12:00pm}\\
  \vspace{0.1in}\large{\textit{\hmwkClassInstructor}} \\
  \vspace{3in}
}

\author{\hmwkAuthorName}
\date{}

% Part Command
\renewcommand{\part}[1]{\textbf{\large Part \Alph{partCounter}}\stepcounter{partCounter}\\}

% Mathematical Commands
% Algorithms
\newcommand{\alg}[1]{\textsc{\bfseries \footnotesize #1}}

% Calculus
\newcommand{\deriv}[1]{\frac{\mathrm{d}}{\mathrm{d}x} (#1)}
\newcommand{\pderiv}[2]{\frac{\partial}{\partial #1} (#2)}
\newcommand{\dx}{\mathrm{d}x}

% Probability and Statistics
\newcommand{\Var}{\mathrm{Var}}
\newcommand{\Cov}{\mathrm{Cov}}
\newcommand{\Bias}{\mathrm{Bias}}
\newcommand*{\prob}{\mathds{P}}
\newcommand*{\E}{\mathds{E}}

% Number Sets
\newcommand*{\Z}{\mathbb{Z}}
\newcommand*{\Q}{\mathbb{Q}}
\newcommand*{\R}{\mathbb{R}}
\newcommand*{\C}{\mathbb{C}}
\newcommand*{\N}{\mathbb{N}}

\begin{document}

\maketitle
\pagebreak

% To change sections, use: \renewcommand{\hmwkSection}{B}
% Example: Uncomment the line below when you start Section B
\renewcommand{\hmwkSection}{B}

\begin{homeworkProblem}
  \begin{enumerate}[(i)]
    \item Show that if $\nabla f(x)$ is the gradient of the quadratic function $f(x) = g^T x + \frac{1}{2} x^T B x$, then for any $s$,
    \[
    Bs = \gamma, \text{ where } \gamma = \nabla f(x + s) - \nabla f(x).
    \]
    \textit{Hint: See Lecture slides Lecture 7.}

    \begin{proof}
      Since $\nabla f(x) = Bx + g$,
      \[
        \gamma = \nabla f(x + s) - \nabla f(x) = Bs.
      \]
    \end{proof}

    \item Consider a general nonlinear function $f : \mathbb{R}^n \to \mathbb{R}$, and apply a quasi-Newton generic linesearch method (GLM) for its minimization that generates iterates $x^k$. Suppose that $B^k$ is a given symmetric matrix (that approximates the Hessian of $f$ on the $k$th iteration), and that $B^{k+1} = B^k + \beta v v^T$ is a rank-one correction for which the secant condition
    \begin{equation}
    B^{k+1} \delta^k = \gamma^k, \text{ where } \gamma^k = \nabla f(x^{k+1}) - \nabla f(x^k) \text{ and } \delta^k = x^{k+1} - x^k, \tag{1}
    \end{equation}
    is required to hold. Show that the resulting $B^{k+1}$ is unique, and give its form; state any assumptions you need to make.

    \begin{proof}
      Note that 
      \[
        \gamma^k = (B^k + \beta v v^T)\delta^k \implies \beta(v^T\delta^k)v = \gamma^k - B^k\delta^k.
      \]
      Write scalar $\rho = (\beta(v^T\delta^k))^{-1}$, so $v = \rho(\gamma^k - B^k\delta^k)$. Assume that $(\gamma^k - B^k\delta^k)^T\delta^k \neq 0$. Then 
      \[
        \gamma^k - B^k\delta^k = \beta\rho^2[(\gamma^k - B^k\delta^k)^T\delta^k](\gamma^k - B^k\delta^k) \implies \beta\rho^2 = \frac{1}{(\gamma^k - B^k\delta^k)^T\delta^k}.
      \]
      Thus,
      \[
        \beta vv^T = \frac{(\gamma^k - B^k\delta^k)(\gamma^k - B^k\delta^k)^T}{(\gamma^k - B^k\delta^k)^T\delta^k},
      \]
      so $B^{k+1}$ is unique.
    \end{proof}

    \item Now suppose that (in the same set-up as in (ii))
    \begin{equation}
    B^{k+1} = B^k + \theta \gamma^k (\gamma^k)^T + \beta v v^T \tag{2}
    \end{equation}
    is a rank-two correction to a given symmetric matrix $B^k$ for which (1) holds. Find an expression for $\theta, \beta$ and $v$ in terms of
    \[
    \phi = 1 - \theta (\gamma^k)^T \delta^k,
    \]
    where $\delta^k = x^{k+1} - x^k$; once again state any assumptions you require. What special cases occur in the formula you have derived when $\phi = 0$ and $\phi = 1$?
    \begin{proof}
      We can easily see that
      \[
        \theta = \frac{1 - \phi}{(\gamma^k)^T\delta^k}.
      \]
      Note that
      \[
        \gamma^k = (B^k + \theta \gamma^k (\gamma^k)^T + \beta v v^T)\delta^k \implies \beta(v^T\delta^k)v = \phi\gamma^k - B^k\delta^k,
      \]
      so $v$ is parallel to $\phi\gamma^k - B^k\delta^k$. Suppose $v = \phi\gamma^k - B^k\delta^k$. Assume that $(\phi\gamma^k - B^k\delta^k)^T\delta^k \neq 0$. Then
      \[
        \beta(v^T\delta^k) = 1 \implies \beta = \frac{1}{(\phi\gamma^k - B^k\delta^k)^T\delta^k}.
      \]
      When $\delta = 1$, then $\theta = 1$ and the update is the same as in (ii). 
    \end{proof}

    \item (difficult) Show that if $B$ is symmetric positive definite and $\delta^T \gamma > 0$, then the BFGS update
    \[
    B^+ = B + \frac{\gamma \gamma^T}{\gamma^T \delta} - \frac{B \delta \delta^T B}{\delta^T B \delta}
    \]
    is also positive definite.
  \end{enumerate}
\end{homeworkProblem}

\newpage

\begin{homeworkProblem}
  Consider the overdetermined system $r : \mathbb{R} \to \mathbb{R}^2, r(x) := (r_1(x) \enspace r_2(x))^T = 0$, with $x \in \mathbb{R}$, where
  \[
  r(x) := (x + 1 \quad \lambda x^2 + x - 1)^T, \quad \text{where } \lambda \in \mathbb{R} \text{ is a constant.}
  \]
  Consider the least-squares objective $f(x) = \frac{1}{2} r(x)^T r(x)$; compute its gradient and Hessian, and show that $x^* = 0$ is a stationary point of $f$ for all $\lambda$, and it is a minimizer of $f$ if $\lambda < 1$.
  Apply Newton's method and Gauss-Newton method, without linesearch (i.e., with stepsize $\alpha^k = 1$, for all $k$) to $f(x)$: compute the Newton iterate in the form $x^{k+1} = \psi_N(x^k)$ and the Gauss-Newton, $x^{k+1} = \psi_{GN}(x^k)$.
  Consider now the functions $\psi_N(x)$ and $\psi_{GN}(x)$. Show that $\psi'_N(0) = 0$ for all $\lambda$, and show that this implies that the rate of convergence of Newton is at least quadratic (asymptotically). Show that $\psi'_{GN}(0) = \lambda$ and that this implies that convergence of Gauss-Newton is linear if $|\lambda| < 1$, and does not occur if $|\lambda| > 1$.

  \textit{Comment: Note that $r(x^*) = r(0) = [1, -1]^T \neq 0$, and so we have a nonzero residual least-squares problem (and hence, the Gauss-Newton approximation to the Hessian is poor, yielding only linear convergence, or even nonconvergence of Gauss-Newton method).}
\end{homeworkProblem}

\end{document}